\pdfoutput=1
\documentclass[11pt]{article}
\usepackage[T1]{fontenc}
\usepackage{lmodern}
\usepackage[margin=.9in]{geometry}
\usepackage{amsmath,amssymb,amsthm,mathtools,microtype}
\usepackage[hidelinks,pdftitle={The canonical polarization phase space of gravitational and Yang--Mills wave packets}]{hyperref}
\newtheorem{theorem}{Theorem}
\newtheorem{corollary}[theorem]{Corollary}
\newcommand{\dd}{\mathrm d}
\newcommand{\eps}{\varepsilon}
\newcommand{\im}{\operatorname{Im}}
\newcommand{\re}{\operatorname{Re}}
\newcommand{\tr}{\operatorname{tr}}
\newcommand{\Ad}{\operatorname{Ad}}
\newcommand{\Div}{\operatorname{div}}
\newcommand{\TF}{\mathrm{TF}}
\newcommand{\TT}{\mathrm{TT}}
\title{The canonical polarization phase space of gravitational and Yang--Mills wave packets}
\author{}
\date{Paper 74}
\begin{document}
\maketitle
\begin{abstract}
We derive a high-frequency limit of the Einstein--Yang--Mills
symplectic form on exact constrained linear wave packets. The
limit is the imaginary part of the Hermitian ray flux of the
joint graviton and color polarization bundle. It gives a
projective two-form at fixed intensity and a matrix Poisson
bracket on the unrestricted packet amplitudes. On two exact
non-Abelian sandwiches a common circular graviton launch gives
the same gravitational polarization covariance at each cut
after integration over a fixed radial aperture. The backgrounds
nevertheless give opposite nonzero canonical responses between
two gravitational polarimeters. A compact pure-gravity
quadrature packet at the first cut reads that response in the
late tidal signal. A separate phase-coherent comparison also
distinguishes the histories. The background metric and all
local gauge-curvature jets agree exactly; the wave comparison
is on exact linearized solutions in their controlled optical
limit.
\end{abstract}

\section{Main statements}
Paper 29B \cite{p29b} constructs oriented polarization measures
for high-frequency vacuum gravity. Paper 29C \cite{p29c}
obtains a rescaled Einstein symplectic pairing on marked cuts of
exact two-stream vacuum families. Paper 60 \cite{p60} proves
physical graviton--color transport and shows that relative
color transport is visible in certain outgoing gravitational
amplitudes and powers. We ask what the coupled propagation does
to the canonical data, and what can be distinguished when all
single-cut gravitational polarization matrices of a specified
aperture are unchanged.

Use signature \((-+++)\), curvature \(F=\dd A+A\wedge A\), and
the action
\[
 S[g,A]=\frac1{2\kappa}\int R\,\dd\mathrm{vol}_g
        -\frac14\int\langle F_{ab},F^{ab}\rangle\,\dd\mathrm{vol}_g,
 \qquad \kappa=8\pi G_N.
\]
The color metric is positive invariant. We linearize only the
wave perturbations on the smooth exact backgrounds in
Section~\ref{sec:bg}. Let \(\phi\) be a fixed oriented eikonal
with \(k=(\dd\phi)^\sharp\) future null, and \(\eps\) its
reduced wavelength (inverse angular frequency in this phase
convention). In the real wave convention
\(\eps\re(e^{i\phi/\eps}(h,a))\) write
\[
 v=(H,a_\mathcal S), \qquad
 H=h_\mathcal S^\TF/(2\sqrt\kappa),\quad
 \mathcal S=k^\perp/\mathbb R k,\quad
 \mathcal V=\mathrm{Sym}^2_0\mathcal S^*
        \oplus(\mathcal S^*\otimes\operatorname{ad}P).
\]
The complex structure records two real optical quadratures. For
\(\mathrm{SU}(2)\) there are eight complex physical channels.

\begin{theorem}[Canonical optical limit]\label{thm:can}
For arbitrary smooth compact mixed profiles \(v_0,w_0\) in a
flat incoming aperture there are real-linear preparations
\(T_\eps\) by exact constrained linear Einstein--Yang--Mills
solutions on a common buffered laboratory. Let \(v,w\) be their
physical leading amplitudes on a pre-caustic Cauchy section
\(\Sigma\) with the packets inside its buffer. With \(n\) the
future normal and the coordinate--momentum convention for
the action symplectic form,
\begin{equation}\label{eq:can}
 \lim_{\eps\downarrow0}\eps^{-1}\Omega_{\mathrm{EYM},\Sigma}
      (T_\eps v_0,T_\eps w_0)
 =\im\int_\Sigma v^\dagger w\,\dd\mu_\Sigma,
 \qquad \dd\mu_\Sigma=-(n\cdot k)\,\dd\Sigma.
\end{equation}
The limit is preserved by the full coupled transport. At fixed
intensity \(I=v^\dagger v>0\), quotienting the common phase gives
\(\tfrac I2\dd\im(z^\dagger\dd z)\), \(z=v/\sqrt I\), per
unit ray flux. For a smooth bounded Hermitian ray matrix \(B\)
write \(F_B(v)=\tfrac12\int v^\dagger Bv\,\dd\mu\). On a common
reference cut the quadratic brackets are
\begin{equation}\label{eq:poisson}
 \{F_B,F_C\}=\frac1{2i}\int\tr(vv^\dagger[B,C])\,\dd\mu.
\end{equation}
Their exact finite-dimensional packet implementation is
specified in Section~\ref{sec:bracket}. No unrestricted
field-theoretic inverse limit is asserted.
\end{theorem}

\begin{theorem}[Same single-cut covariance, different history]\label{thm:benchmark}
Two exact SU(2) sandwiches have the same metric and Yang--Mills
curvature, including all local covariant curvature derivatives.
Their gap transports come from reordering the same three color
pulses. For a common circular graviton-only launch and a
sufficiently small fixed radial aperture, the entire optical
\(2\times2\) gravitational polarization covariance integrated
on each cut agrees. Nevertheless the signed linear
gravity-polarization power on two separated cuts has a
canonical bracket of equal magnitude and opposite nonzero
signs. This is the bracket on the full physical packet space
evaluated at the graviton-only launch. A compact pure-gravity
quadrature packet at the first cut produces the corresponding
late tidal response. A separate passive comparison of the
imaginary same-ray mutual coherence between both cuts also
has opposite nonzero signs.
\end{theorem}

The equal single-cut matrices involve integration over
transverse ray labels, not coherent interference across
different labels, and hold at leading optical order. Equality
of the full resolved field, color state, or arbitrary
finite-wavelength measurements is not claimed.

\section{Transport and exact incident packets}\label{sec:prep}
A compensated infinitesimal gauge is
\(K(\zeta,\chi)=(\mathcal L_\zeta g,D\chi+\iota_\zeta F)\).
The on-shell Hessian obeys \(PK=0=K^\dagger P\).
Action-adjoint de Donder and Lorenz gauges give scalar wave
principal parts in the field and subsidiary gauge-violation
equations, with the kinetic row normalizations
\cite{hack,p60}. Thus zero gauge data and the linearized
constraints yield exact physical solutions. At null \(k\)
the de Donder and Lorenz symbols are
\(k^a(h_{ab}-g_{ab}\tr h/2)=0\) and \(k^a a_a=0\).
Metric and color principal gauge additions do not affect their
screen readings. The compensated term \(\iota_\zeta F\) is
one order lower for a gauge producing a wave-scale metric
perturbation.

In a parallel screen frame put \(f_A=F(k,e_A)\), and write
\(C_f a=(\langle f_{(A},a_{B)}\rangle)^\TF\) and
\((C_f^*H)_A=H_{AB}f_B\). Paper 60 gives
\begin{equation}\label{eq:transport}
  (\nabla_k^{\mathcal V}+\tfrac12\Div k)v
   =\mathcal A_fv,\qquad
  \mathcal A_f=\sqrt\kappa
       \begin{pmatrix}0&-C_f\\ C_f^*&0\end{pmatrix}.
\end{equation}
The factors follow from the stress variation \(2iC_fa\)
at the fast derivative and the derivative of the Yang--Mills
density. In the latter, the metric trace term and the first
raised-index variation cancel by the de Donder symbol;
the remaining screen insertion is \(-i C_f^*h^\TF\).
The resulting unscaled source terms are
\(-2\kappa C_fa\) and \(\tfrac12 C_f^*h^\TF\).
This yields the displayed skew-adjoint mixing.
The color slot uses the actual adjoint connection.
An action-current limit additionally needs the configuration
errors to be smaller than the wave amplitude.

\subsection{Local exact constraints}
Use a larger spacelike slice with only a flat incoming
subpatch required to contain the perturbation support.
All data vanish on the rest of the slice, including its
curved part. Choose the laboratory and its time and spatial
buffers so that the past domains needed by the receiver
patches lie inside this larger slice. A slice confined
entirely to the flat support subpatch would not suffice.
In flat coordinates define
\[
 Q_{ij}=-\Delta\delta_{ij}+\partial_i\partial_j,\qquad
 \mathbb T_{ij}[b]=Q_i{}^kQ_j{}^\ell b_{k\ell}
                   -\tfrac12 Q_{ij}Q^{k\ell}b_{k\ell}.
\]
The first operator is identically divergence-free. The
second is both divergence-free and tracefree by
\(Q^2=(-\Delta)Q\) and \(\tr Q=-2\Delta\).
For a unit constant spatial phase covector apply
\(\eps^3 Q\) to the compact connection carrier seed \(a_0\)
and \(\eps^2Q\) to its \(-ia_0\) velocity seed. Apply
\(\eps^5\mathbb T,\eps^4\mathbb T\) to the metric seeds
\(h_0=2\sqrt\kappa H_0,-ih_0\), respectively.
Take real parts. Their leading symbols on the desired
transverse seeds give the exact requested optical amplitudes.
For a nonconstant nonzero spatial phase gradient on a flat
subpatch, place \(|\dd_\Sigma\phi|^{-2}\) or
\(|\dd_\Sigma\phi|^{-4}\) in the respective seed and include
the known local normal angular frequency in the velocity
seed. All exact corrections retain compact support.
The temporal components and their gauge velocities vanish
in this flat chart. Here \(A=F=0\), so the linearized
stress does not spoil the vacuum TT constraints, and
Gauss is ordinary transverse electric data. Zero extension
within the larger slice preserves these exact constraints.
The subsidiary equation gives exact constrained evolution.

\subsection{Energy estimate at the required scale}
An \(O(\eps)\) wave-energy error alone would not give the
symplectic limit, since the packet configuration is itself
\(O(\eps)\). The compact differential preparations have
finite expansions with smooth coefficients. Continue their
initial phase through the two branches \(\phi_+,\phi_-\)
with opposite initial normal derivatives, \(\phi_+=\phi\).
In the flat coordinates \(t=(u+v)/\sqrt2\),
\(z=(v-u)/\sqrt2\) at \(t=0\), the phases can start with
\(-v\) and \(u\) (up to a common scale). The second phase
extends exactly as a null phase. For a later flat preparation
with a smoothly varying gradient, a common pre-caustic
laboratory for both branches is available by taking weak
focusing and a sufficiently short fixed geometry.
The intervening color rotations do not affect the eikonals.
Use the complex approximate solution
\[
 U_{\rm app}=\eps e^{i\phi_+/\eps}U_0^+
            +\eps^2 e^{i\phi_+/\eps}U_1^+
            +\eps^2 e^{i\phi_-/\eps}U_1^- .
\]
Solve full leading transport for \(U_0^+\), including its
longitudinal representatives. Solve a sourced transport
for \(U_1^+\) to cancel the order-\(\eps\) residual, and a
homogeneous one for \(U_1^-\). Their initial sum matches
the order-\(\eps^2\) configuration. Their weighted
difference with opposite normal frequencies, after the
known slow derivative of \(U_0^+\) is subtracted, matches
the order-\(\eps\) velocity. This two-by-two system is
invertible. The residual is \(O(\eps^2)\) in \(L^2\) and
remains so after one \(\eps\)-weighted derivative.
The initial mismatch is \(O(\eps^3)\) in configuration
before differentiation, \(O(\eps^2)\) in velocity.
Local hyperbolic energy estimates on the fixed buffered
region give \(O(\eps^2)\) errors in \(H^1\times L^2\)
against \(U_{\rm app}\). The weighted higher estimate
gives \(\eps\nabla^2(U-U_{\rm app})=O_{L^2}(\eps^2)\)
on the chosen cuts: apply the one-higher wave energy
estimate with weight \(\eps\) and smooth background
coefficient bounds. All relevant initial weighted
derivatives of the mismatched polynomial data are of
the same order. Hence in particular
\begin{align}
 U&=\eps\re(e^{i\phi/\eps}U_0^+)+O_{L^2}(\eps^2),\notag\\
 \nabla U&=\re(ie^{i\phi/\eps}\dd\phi\otimes U_0^+)
                         +O_{L^2}(\eps).\label{eq:errors}
\end{align}
Higher weighted control gives the flat tidal-field reading
used below. Bounds are uniform on fixed finite-dimensional
input spaces, by linearity. This is not a differentiated
estimate for a family of exact nonlinear backgrounds.

\section{Normalizing the action current}
The action current \cite{leewald} on a spacelike cut is
the gravitation and Yang--Mills configuration--momentum
wedge, modulo a spatial boundary divergence. The compact
buffer removes that divergence. Take
\(\Omega_{\mathrm{EYM}}=-\dd_{\rm fields}\Theta\), with
\(\Theta=\int(\pi^{ij}\delta q_{ij}+\langle E^i,\delta A_i\rangle)\).
An exact compactly supported linear gauge change removes
the \emph{leading} temporal components near a cut: its
oscillatory parameters are of size \(\eps^2\), and inversion
of the principal temporal gauge map uses \(n\cdot k\ne0\).
The background compensation in the color slot is then
subleading. This change has zero edge symplectic charge.
In this representative the leading gauge conditions imply
spatial TT metric and transverse color.

For \(K=\tfrac12\mathcal L_nq\), the momenta are
\(\pi^{ij}=\sqrt{\det q}(K^{ij}-\tr K\,q^{ij})/(2\kappa)\)
and \(E^i=\sqrt{\det q}\,q^{ij}F_{nj}\).
Their leading fast variations are
\[
 \delta\pi^{ij}
   =\frac{\sqrt{\det q}}{4\kappa}\partial_n h_{\rm act}^{ij,\TT}
                  +O_{L^2}(\eps),\qquad
 \delta E^i
   =\sqrt{\det q}\,q^{ij}D_n a_{j,\rm act}
                  +O_{L^2}(\eps),
\]
where the subscript indicates actual perturbations,
including their factor of \(\eps\). Background-density,
lapse and shift terms at this order are lower, and
the YM kinetic momentum has no fast metric derivative.
The tensor coefficient becomes exactly the color
coefficient after \(h^\TT/(2\sqrt\kappa)\).
Configuration errors in \eqref{eq:errors} pair with a
derivative of order one, and derivative errors pair with a
configuration of order \(\eps\), giving only \(O(\eps^2)\).
At order \(\eps\), use the pointwise real-part identity
\[
 (\re(e^{i\vartheta}v),\re(ie^{i\vartheta}w))_{\mathbb R}
 -(\re(e^{i\vartheta}w),\re(ie^{i\vartheta}v))_{\mathbb R}
 =-\im(v^\dagger w).
\]
The double carrier cancels before integration.
Multiplication by \(\eps\,n\cdot k\) proves \eqref{eq:can}.
Equation~\eqref{eq:transport} gives
\(k(v^\dagger w)+(\Div k)v^\dagger w=0\), so the entire
Hermitian flux is preserved. In common ray labels after
division by real optical spreading the transfer is unitary.
The separate gravity and color symplectic fluxes need
not be conserved.

For \(\Omega(v,w)=\im\int v^\dagger w\,\dd\mu\), the
convention \(\iota_{X_B}\Omega=\dd F_B\) gives
\(X_B=-iBv\) and \(\{F_B,F_C\}=\Omega(X_B,X_C)\)
as in \eqref{eq:poisson}. At fixed intensity, writing
\(v=\sqrt I\,z\) makes the common amplitude phase radical,
and \( \dd\im(z^\dagger\dd z)(\delta z,\delta' z)
=2\im(\delta z^\dagger\delta' z)\).
The quotient is the usual projective form \cite{ashtekar}.
It gives a \(\mathbb{CP}^7\) joint fiber; the gravity
\(\mathbb{CP}^1\) subfiber at flat entry is not generally
preserved by the slabs. This phase quotient is not
a finite-wavelength gauge symmetry. On the gravity-only
one-carrier restriction \(v=Fz/(2\sqrt\kappa)\) for strain
amplitude \(F\), its coefficient is
\(F^2/(8\kappa)=F^2/(64\pi G_N)\) against the incident
affine area measure, the same positive normalization
as the marked incoming cut of \cite{p29c}. This does
not import that paper's nonlinear vacuum limiting
argument into mixed EYM. Paper 29B's rank-one
oriented covariance fixes the projective ray; its trace
does not fix polarization.

\section{Two orderings with the same optical marginals}\label{sec:bg}

Take orthonormal adjoint generators \(T_a\) with
\([T_a,T_b]=\epsilon_{abc}T_c\), two positive smooth
disjoint active pulses \(f_1,f_2\), and \(f=f_1+f_2\).
Consider the Brinkmann sandwiches
\begin{align}
 g&=-2\dd u\,\dd v+(\dd x^1)^2+(\dd x^2)^2
                     -\kappa f(u)^2|x|^2(\dd u)^2,\notag\\
 A_\pm&=\big[C_\pm(u)-f(u)(x^1T_1+x^2T_2)\big]\,\dd u,
 \label{eq:background}
\end{align}
with \(C_\pm\) smooth and supported strictly inside the
gap between the active pulses. This is an exact
Coleman--G\"uven family \cite{coleman,guven,gst,p60}.
In the displayed gauge \(F_{uA}=fT_A\) are the only
independent curvature entries. The Yang--Mills equation
follows directly because only \(F^{vA}\) and its
antisymmetric partners are nonzero, \(\det g=-1\),
and the field is independent of \(v\) with \(A_v=A_A=0\).
The only stress entry is \(T_{uu}=2f^2\), and
\(\operatorname{Ric}_{uu}=2\kappa f^2\), by the transverse
Laplacian of the profile. Thus they solve the full coupled
background equations. Where \(C_+\ne C_-\), the curvature
vanishes on a neighborhood; elsewhere the background
connections coincide where needed. Their curvature tensors
and all local covariant curvature derivatives agree in
the indicated gauge.

Use the adjoint color rotations \(R_a(\alpha)\) about
\(T_a\). Choose \(C_\pm\) in three disjoint intervals so
that its actual parallel action in the gap is
\begin{equation}\label{eq:word}
 W_\pm=
   R_1(\pm\pi/2)R_2(\theta)R_1(\mp\pi/2)=R_3(\pm\theta),
      \qquad 0<\theta<\pi .
\end{equation}
The products are operators, with the rightmost rotation
first in time. Both use the same three signed rotation
pulses, changing their order. A bump along a fixed generator
has adjoint rotation given by its integral with the sign
from the parallel equation. The ordered noncommuting active
pair \(fT_1,fT_2\) fixes the adjoint endpoint frames.
A gauge transformation intertwining the displayed identical
curvatures on both active slabs must act trivially in their
adjoint frames; it cannot turn one relative transporter into
the other. A gap-only potential has no absolute gauge
meaning. In particular it is possible to gauge away an
unmarked \(u\)-only term at the cost of rotating the later
curvature data \cite[Secs.~II and IV]{gst}; that does not
preserve the two marked endpoint frames of this comparison.

Use the counterpropagating planar incident congruence with
\(k^u=1\) as its affine scale. A common constant
rescaling of its optical covector can be absorbed in
\(\eps\) and the flux normalization. The two transverse
ray labels \(x_0\) at flat entry obey
\[
 x(u)=Y(u)x_0,\qquad Y''+\kappa f(u)^2Y=0,\qquad
 Y=1,\quad Y'=0\quad\text{at entry}.
\]
The small active pulses and bounded interval can be chosen
so \(Y>0\) everywhere: e.g., if the largest Volterra
integral of \((u-s)\kappa f(s)^2\) is \(\delta<1/2\),
then \(\|Y-1\|_\infty\le\delta/(1-\delta)<1\).
The parallel screen frame can be represented by
\(e_A^u=0,e_A^i=\delta_A^i\), with \(e_A^v\) set by
\(g(e_A,k)=0\). In particular the screen colors are \(fT_A\)
even at noncentral rays. The active adjoint connection
along a ray is \(-fY(u)x_0^AT_A\); it need not vanish
off-axis and is retained throughout. Smooth amplitudes
below are in common ray coordinates with the real
geometrical spreading factored out. The longitudinal
packet factor and its ray-label measure are compact and
common; they do not affect the rotational comparison.

For \(O\in O(2)\), simultaneous transverse reflection or
rotation \(O\), screen transformation \(O\), and color
rotation \(\widehat O=\operatorname{diag}(O,\det O)\)
are exact symmetries of the active propagation, including
the off-axis connection. On either entire gap the gravity
propagates freely, and the color action is the same
\(W_\pm\) for every label, on each screen slot. Completed
gap transport \(W_\pm=R_3(\pm\theta)\) commutes with every
proper \(\widehat O\) and is exchanged between the two
backgrounds by
\(\widehat O_*=\operatorname{diag}(1,-1,-1)\),
where \(O_*=\operatorname{diag}(1,-1)\).
Thus the transfer of the gravity amplitude to \emph{each}
common optical cut has proper-rotation covariance and a
reflection that exchanges the two experiments. At cuts
inside the gap this requires only gravity's free transport.
No such symmetry of each partial gap \emph{color} history
is assumed. The same argument works for rotation-invariant
spacelike sections of the tube. Take the early cut \(E\)
and the late cut \(L\) in inertial flat spacelike sections
strictly after slab 1 but before any gap pulses, and
strictly after slab 2, respectively, on the receiver
supports. Their cuts and apertures can be chosen in these
open common flat buffers.

\subsection{What the aperture forgets}

In the screen choose
\(H_+=\operatorname{diag}(1,-1)/\sqrt2\) and
\(H_\times=(\begin{smallmatrix}0&1\\1&0\end{smallmatrix})/\sqrt2\).
Let \(z_\circ=(1,i)^t/\sqrt2\) in this pair, and launch
the same \(b(|x_0|)z_\circ\) with zero incoming color,
for \(b\) smooth real, compact and positive around the
central ray. Let \(H_\pm(s,x_0)\) be the gravitational
amplitudes on a matched cut, and \(d\mu\) the common
positive ray-label measure including a smooth compact
radial aperture weight. For a proper screen rotation
the induced real spin-two matrix \(D(O)\) fixes the
circular input up to one overall phase. Proper-rotation
covariance implies
\[
 \mathcal D_\pm(s)
   =\int H_\pm(s,x_0)H_\pm(s,x_0)^\dagger\,\dd\mu
                 \in\operatorname{span}_{\mathbb R}(I,\sigma_y),
 \qquad \sigma_y=\begin{pmatrix}0&-i\\i&0\end{pmatrix}.
\]
Indeed the integrated matrix commutes with the real
spin-two rotations, and its Hermitian part of that
commutant is precisely the displayed span. The
\(O_*\) reflection acts on the gravitational pair
by \(P=\operatorname{diag}(1,-1)\), and
\(P\overline{z_\circ}=z_\circ\). Its composition with
complex conjugation of the \emph{amplitudes}, without
reversing the eikonal, intertwines the two real physical
transports in the optical calculation. Therefore
\(H_-(s,O_*x_0)=P\overline{H_+(s,x_0)}\).
It follows that
\(\mathcal D_-(s)=P\overline{\mathcal D_+(s)}P^t
=\mathcal D_+(s)\), since \(P\sigma_yP^t=-\sigma_y\).
This is equality of all four entries of the matrix
on every such cut, not just of its trace or of a
particular linear analyzer. It is not a pointwise
off-axis equality of all polarization components.

\section{Canonical history and its measured response}\label{sec:bracket}

Write
\[
 \Theta_i=\sqrt\kappa\int f_i(u)\,\dd u,\quad
 c_i=\cos\Theta_i,\quad s_i=\sin\Theta_i,\qquad
 K_0=\begin{pmatrix}0&-1\\1&0\end{pmatrix}.
\]
Use small nonzero \(\Theta_i\), so all relevant \(c_i,s_i\)
are positive. At the central ray the color injection
\(J=C^*_{(T_1,T_2)}\) is an isometry with columns
\((T_1,-T_2)/\sqrt2,(T_2,T_1)/\sqrt2\).
As \(A(k)=0\) on this ray in each active slab,
\eqref{eq:transport} reduces there to a rotation
by \(\Theta_i\) mixing \(H\) and \(JH\), acting
trivially on the orthogonal color complement.
An incident unit \(z_\circ\) gives
\(v_E=(c_1z_\circ,s_1Jz_\circ)\) at the early cut.
The map from all eight early channels to late
gravity is
\[
 L_\pm=\big(c_2 I,\ -s_2J^\dagger(W_\pm\oplus W_\pm)\big)
\]
on the central ray. Because
\(J^\dagger(W_\pm\oplus W_\pm)J
      =\cos\theta\,I\pm\sin\theta\,K_0\) and
\(K_0z_\circ=-iz_\circ\), the late central gravity is
\begin{equation}\label{eq:late}
 H_{L,\pm}=(c_1c_2-s_1s_2\cos\theta
                       \ \pm i s_1s_2\sin\theta)\,z_\circ.
\end{equation}
In particular \emph{even the central output ray} agrees.
The finite-aperture proof above is needed for all of the
equal marginal measurements.

Use the same real signed linear polarimeter
\(Z=\operatorname{diag}(1,-1)\) at \(E\) and \(L\).
It reads \(F_Z\), half the difference of the two linear
polarization powers, calibrated in physical units of
\(H\). On the common early full packet space the two
effects are \(B_E=\operatorname{diag}(Z,0)\) and
\(B_{L,\pm}=L_\pm^\dagger Z L_\pm\); away from the axis
use the actual smooth physical transport followed by
the same gravitational compression. The symbols can
include smooth matched aperture weights. At the axis
the gravity part of \(L_\pm^\dagger\) is \(c_2I\);
since \(Z^2=I\),
\[
 v_E^\dagger B_E B_{L,\pm} v_E
        =c_1c_2\, z_\circ^\dagger H_{L,\pm}.
\]
It follows from \eqref{eq:poisson} that their brackets are
\begin{equation}\label{eq:witness}
 \{F_{Z,E},F_{Z,L}\}_{\pm}
       =\ \pm c_1c_2s_1s_2\sin\theta
\end{equation}
per unit incident intensity on the central ray. This is
the full bracket, not an inverse on a gravity-only
subspace. A reflection-symmetric finite aperture retains
the strict signs by continuity on a small disk.
Moreover the reflecting conjugation of the \emph{full}
early-to-late transfer is antiunitary on marked rays,
so it reverses \eqref{eq:can}; it preserves both real
\(Z\) effects under the choice of matched weights.
Hence the integrated brackets have exactly equal magnitude
and opposite signs. Pointwise equality of two passive
values of \(F_Z\) at a chosen input would not determine
these derivatives on the incident canonical directions.

\subsection{Exact finite-dimensional bracket limit}

For any finite set of compact optical modes spanning a
complex space \(M\), the exact prepared solution maps
pull \(\Omega_{\rm EYM}/\eps\) back to matrices converging
to the nondegenerate restriction of \eqref{eq:can}.
Their inverses converge in this fixed dimension.
One can choose \(M\) at the benchmark to contain the
input and both optical Hamiltonian directions (pulling
them back to the flat launch labels), for both backgrounds.
The restricted limiting brackets at that point are the
full optical ones; the vectors \(-iB_E v\) and
\(-iB_Lv\) already lie in the chosen space.

In the common flat receiver regions the linearized
tidal Weyl tensor is diffeomorphism gauge invariant.
For the scaled field \(\eps\,\delta W(n,e_A,n,e_B)^\TF\),
the real wave factor at leading order is
\((n\cdot k)^2h^\TF/2\), up to a fixed Riemann-sign
convention. Divide by this known factor between \(h^\TF\)
and \(H\), and by geometrical spreading where applicable.
The signed integrated square of these calibrated
\emph{real} tidal components converges to \(F_Z\):
each square gives \(1/2\) by averaging, and the spatial
covector of a nonzero null phase on a spacelike cut
never vanishes, so the \(2\phi/\eps\) terms vanish
against smooth compact cut weights. The weighted
second-derivative estimate of Section~\ref{sec:prep}
controls the errors. By linearity on \(M\), these
quadratic readouts converge coefficientwise and in \(C^1\)
on bounded sets. Their brackets with the exact pulled-back
normalized symplectic matrices converge to the indicated
optical brackets. Equivalently one multiplies the
brackets computed with the \emph{unscaled} symplectic
form by \(\eps\). This is a finite-mode result, not a
limit of arbitrary localized Peierls brackets.

\subsection{A purely gravitational quadrature drive}

There is a direct response experiment distinct from
taking an inverse of a large set of data. On the early
inertial spacelike section in the common flat buffer,
superpose an independently prepared small optical
packet with leading variation
\[
 X_{F_{Z,E}}(v_E)=-i(\chi_E ZH_E,0),
\]
where \(\chi_E\) is the compact early aperture. Its
polarization support lies within the shared \(F=A=0\)
subpatch. The exact local TT differential construction
of Section~\ref{sec:prep}, now including the
\(|\dd_\Sigma\phi|^{-4}\) symbol factor and the known
normal angular frequency, implements this compact
quadrature direction with zero color data. Its larger
intermediate Cauchy slice covers the late support.
Linearity allows addition of this exact constrained
solution to the baseline after the early cut. The same
drive works in the two comparisons because the
pre-gap \(H_E\) is common. The late calibrated
tidal intensity derivative is
\[
 \left.\partial_\eta
       F_{Z,L}(v_E+\eta X_{F_{Z,E}}(v_E))\right|_{\eta=0}
       =-\{F_{Z,E},F_{Z,L}\}(v_E)
\]
in the stated convention. Its optical limit is realized
by the exact packet and detector estimates.
This drive is a new independent preparation at \(E\);
one does not abruptly change a source-free solution's
past. If a compact time-dependent linear source is
desired, evolve the auxiliary TT data a short time
into the early flat buffer and multiply its free
solution by a smooth time ramp from zero to one.
The linear forced Einstein tensor in that buffer is
then supported in the ramp and is proportional to
the combination of its TT field and first time
derivative with the ramp's first two derivatives.
Its temporal components are zero and its spatial
divergence vanishes identically, so the compact
source is conserved and produces the same
retarded auxiliary packet. This is an externally
prescribed linear source, not a nonlinear matter
apparatus or an assertion of positive energy for
the prescribed source.

\subsection{An independent passive comparison}

A synchronized tidal measurement can instead compare
the complex same-label coherence
\[
 \mathcal C_\pm(E,L)=\int
   \left\langle \mathcal P_{L\leftarrow E}H_\pm(E,x_0),
                         H_\pm(L,x_0)\right\rangle\,\dd\mu .
\]
Here \(\mathcal P\) is the common geometric parallel
identification of screen frames along each ray.
Reflection and conjugation give
\(\mathcal C_-=\overline{\mathcal C_+}\).
By \eqref{eq:late} the imaginary central densities
are \(\pm |b(0)|^2c_1s_1s_2\sin\theta\), so they
give opposite nonzero signs on a small fixed aperture.
With two flat tidal receivers, demodulate
\(2e^{-i\phi/\eps}\) times the scaled real Weyl signal
and remove the conjugate branch using a local
frequency filter: its phase \(-2\phi/\eps\) is
separated by the nonzero spatial wave covector.
The error estimates give convergence of these complex
readings in \(L^2\); their bounded aperture product
then converges to the displayed coherence.
A common raywise rephasing cancels, and a change
of either screen frame also changes its parallel
identification, leaving the result fixed. Independent
rephasings at the two cuts are not allowed in this
common-reference comparison. This passive
measurement is not the definition of the canonical
bracket. It uses two coherent readings instead of the
calibrated drive and late signed intensity alone.

\section{Relation to earlier optical and symplectic results}

The covariant action phase space, projective two-form
and quadratic moment-map bracket are established
structures \cite{leewald,ashtekar}. Weak short-wave
gravitational--electromagnetic conversion and conserved
sum currents were studied before \cite{tokuoka};
graviton and gauge tensor oscillations in a locked
Yang--Mills background are not new \cite{cdm}.
The Berry-type effective ray corrections of
\cite{anderson} concern subleading propagation for
linearized gravity and are not the profile Poisson
limit here. The exact plane-wave stock and earlier
gauge-curvature ambiguity \cite{coleman,guven,gst}
are distinct from this gravity-only history comparison.
Papers 29B, 29C and 60 supply oriented polarization,
a separate vacuum symplectic limit, and a coupled
linear transport law, respectively. The new steps
are the mixed canonical action normalization on exact
constrained packets and the all-marginals-equal
canonical and coherent history test. We do not
construct a nonlinear EYM high-frequency limit,
a quantum state or a quantization.

\paragraph{Provenance.}
This analytical manuscript was prepared with AI assistance
and independently checked within the project. It has not
undergone external expert peer review.

\small
\begin{thebibliography}{12}
\setlength{\itemsep}{2pt}
\bibitem{leewald}
J.~Lee and R.~M.~Wald, \emph{Local symmetries and constraints},
J.\ Math.\ Phys.\ \textbf{31}, 725--743 (1990),
\href{https://doi.org/10.1063/1.528801}{doi:10.1063/1.528801}.
\bibitem{hack}
T.-P.~Hack and A.~Schenkel, \emph{Linear bosonic and fermionic quantum gauge theories on curved spacetimes},
\href{https://arxiv.org/abs/1205.3484}{arXiv:1205.3484}.
\bibitem{ashtekar}
A.~Ashtekar and T.~A.~Schilling, \emph{Geometrical formulation of quantum mechanics},
\href{https://arxiv.org/abs/gr-qc/9706069}{arXiv:gr-qc/9706069}.
\bibitem{tokuoka}
T.~Tokuoka, \emph{Interaction of electromagnetic and gravitational waves in the weak and short wave limit},
Prog.\ Theor.\ Phys.\ \textbf{54}, 1309--1317 (1975),
\href{https://doi.org/10.1143/PTP.54.1309}{doi:10.1143/PTP.54.1309}.
\bibitem{cdm}
R.~R.~Caldwell, C.~Devulder and N.~A.~Maksimova,
\emph{Gravitational wave--gauge field oscillations},
\href{https://arxiv.org/abs/1604.08939}{arXiv:1604.08939}.
\bibitem{anderson}
L.~Andersson, J.~Joudioux, M.~A.~Oancea and A.~Raj,
\emph{Propagation of polarized gravitational waves},
\href{https://arxiv.org/abs/2012.08363}{arXiv:2012.08363}.
\bibitem{coleman}
S.~Coleman, \emph{Non-Abelian plane waves}, Phys.\ Lett.\ B
\textbf{70}, 59--60 (1977),
\href{https://doi.org/10.1016/0370-2693(77)90344-6}{doi}.
\bibitem{guven}
R.~G\"uven, \emph{Solution for gravity coupled to non-Abelian plane waves},
Phys.\ Rev.\ D \textbf{19}, 471 (1979),
\href{https://doi.org/10.1103/PhysRevD.19.471}{doi}.
\bibitem{gst}
M.~G\"urses, T.~\c{C}.~\c{S}i\c{s}man and B.~Tekin,
\emph{Einstein--Yang--Mills plane waves: uniqueness and an exact curvature double copy},
\href{https://arxiv.org/abs/2611.12204}{arXiv:2611.12204}.
\bibitem{p29b}
Research archive, \emph{Physical polarization of high-frequency vacuum gravity},
\href{https://github.com/g3173028019-beep/ai-automated-quantum-gravity-research/tree/main/papers/29b-high-frequency-graviton-polarization}{Paper 29B}.
\bibitem{p29c}
Research archive, \emph{A polarization symplectic form at a fixed two-stream null-dust limit},
\href{https://github.com/g3173028019-beep/ai-automated-quantum-gravity-research/tree/main/papers/29c-high-frequency-polarization-symplectic-form}{Paper 29C}.
\bibitem{p60}
Research archive, \emph{Gravitons distinguish non-Abelian transport at fixed local curvature},
\href{https://github.com/g3173028019-beep/ai-automated-quantum-gravity-research/tree/main/papers/60-graviton-nonabelian-holonomy}{Paper 60}.
\end{thebibliography}
\end{document}
